% !PRINTED 507
\weberpage{507}

For all the permutations in $P$ that leave the figure $0$ unchanged
form a divisor $Q$ of $P$. Moreover, on account of the transitivity,
there are in $P$ permutations $\pi_1, \pi_2 \ldots \pi_{m-1}$ which carry the
figure $0$ into $1, 2 \ldots m-1$, and we have

\[
P = Q + Q\pi_1 + Q\pi_2 + \cdots + Q\pi_{m-1},
\]

hence the degree of $P$ is equal to $m$ times the degree of $Q$.

\begin{center}
§. 155.

\gesperrt{Reduction of the Galois resolvent by adjunction.\\
Normal divisor of a group.}
\end{center}

We have now to consider the gradual reduction of the group
of a given equation, which can come about through the adjunction
of certain algebraic quantities.

Let, as in the preceding section, $P$ be the \gesperrt{Galois}
group of degree $p$ and $Q$ one of its divisors of degree $q$
and of index $j$; further let $\psi$ be a function of the roots
belonging to $Q$, and let

\weq{1}{\psi, \psi_1, \psi_2 \ldots \psi_{j-1}}

be the conjugate quantities. We place the following theorem
at the head.

1. \gesperrt{The conjugate quantities (1) are the roots
of an irreducible equation of degree $j$
in $\Omega$.}

For if $t$ denotes a variable, then

\weq{2}{\varphi(t) = (t - \psi)(t - \psi_1) \ldots (t - \psi_{j-1})}

remains unchanged (by §.154, 5.) when a permutation from $P$ is
applied; hence the coefficients of the function $\varphi(t)$, whose
roots are the quantities (1), are contained in $\Omega$ [§. 149, b)].

To prove the irreducibility of $\varphi(t)$, we assume
that $\Phi(t)$ is any function in $\Omega$ which vanishes
for $t = \psi$, so that $\Phi(\psi) = 0$. Since all the permutations
of $P$ may be applied to this equation, it follows that
$\Phi(\psi_1), \Phi(\psi_2) \ldots \Phi(\psi_{j-1})$ must also be zero, and hence that $\Phi(t)$
must be divisible by $\varphi(t)$. But in this the irreducibility
is contained.
